\documentclass[11pt]{article}

\usepackage[margin=1in]{geometry}
\usepackage{amsmath,amssymb,amsthm,mathtools}
\usepackage{algorithm,algpseudocode}

\newtheorem{theorem}{Theorem}
\newtheorem{lemma}[theorem]{Lemma}
\newtheorem{observation}[theorem]{Observation}
\newtheorem{corollary}[theorem]{Corollary}
\newtheorem{claim}[theorem]{Claim}

\newcommand{\cut}{\delta}
\newcommand{\OPT}{\mathrm{OPT}}
\newcommand{\Etree}{\mathbb{E}_{T\sim\mu}}

\begin{document}

\title{$O(\log n)$-approximation for one-sided quasi-expansion\\
via Räcke's cut-sparsifier decomposition}
\date{}
\maketitle

\section{Problem setup}

Let $G=(V,E,c,w)$ be an undirected graph where $c\colon E\to\mathbb{R}_{>0}$
are edge capacities and $w\colon V\times V\to\mathbb{R}_{\ge 0}$ is a symmetric
demand function. Write
\[
  \pi(u) \;=\; \sum_{v\in V} w(u,v)
\]
for the total demand incident to $u$ and $W_{\mathrm{tot}}=\sum_{u\in V}\pi(u)$.
For a nonempty proper subset $S\subsetneq V$ define
\[
  h(S)\;=\;\frac{c_G(\cut(S))}{w_G(S,V)},
  \qquad
  w_G(S,V)\;=\;\sum_{u\in S}\pi(u),
\]
where $c_G(\cut(S))=\sum_{e\in\cut(S)}c(e)$ is the capacity of the cut.
We wish to compute (or approximate) the \emph{one-sided quasi-expansion}
\[
  \psi \;=\; \min_{\emptyset\neq S\subsetneq V} h(S).
\]

\section{Reduction to an anchor graph}

\subsection{Auxiliary instance}

Introduce a new \emph{anchor} vertex $r\notin V$ and set $V'=V\cup\{r\}$.
Define demand pairs on $V'$ by
\[
  D(u,r)=D(r,u)=\pi(u)\quad(u\in V),
  \qquad
  D(u,v)=0\quad(u,v\in V).
\]
The capacity function is extended by setting $c_{G'}=c_G$ on $E$ and $c_{G'}(e)=0$
for any hypothetical edge incident to $r$ (equivalently, $r$ is an isolated
vertex in the capacity graph).

For a cut $X\subseteq V'$ define the \emph{demand cut} $D(\cut(X))$
and \emph{capacity cut} $c_{G'}(\cut(X))$ in the usual way.

\subsection{The anchor identity}

\begin{lemma}[Anchor identity]\label{lem:anchor}
  Let $X\subseteq V'$ and suppose $r\notin X$ (taking complement if necessary so
  that $X$ is the side not containing $r$). Then
  \[
    D(\cut(X)) \;=\; w_G(X,V),
    \qquad
    c_{G'}(\cut(X)) \;=\; c_G(\cut(X\cap V)).
  \]
\end{lemma}

\begin{proof}
  Write $S=X\cap V$ (note $X=S$ since $r\notin X$, so $X\subseteq V$).
  By definition of $D$, the demand crossing the cut $(X,V'\setminus X)$ counts
  every pair $\{u,v\}$ with $u\in X$ and $v\notin X$.
  Since $D(u,v)=0$ for $u,v\in V$, only pairs involving $r$ contribute:
  \[
    D(\cut(X))
    =\sum_{u\in S} D(u,r)
    =\sum_{u\in S}\pi(u)
    =w_G(S,V).
  \]
  For the capacity, since $r$ has no capacity edges, every capacity edge in
  $\cut_{G'}(X)$ has both endpoints in $V$, giving
  $c_{G'}(\cut(X))=c_G(\cut(S))$.
\end{proof}

\begin{corollary}\label{cor:trivial-cut}
  The cut $\{r\}$ satisfies $D(\cut(\{r\}))=W_{\mathrm{tot}}$
  and $c_{G'}(\cut(\{r\}))=0$.
\end{corollary}

\subsection{Bounded-demand reformulation}

Fix a parameter $K\in(0,W_{\mathrm{tot}})$ and define
\[
  \OPT_K(G')
  \;=\;
  \min_{\substack{X\subseteq V'\\ 0<D(\cut(X))\le K}}
  \frac{c_{G'}(\cut(X))}{D(\cut(X))}.
\]
By Lemma~\ref{lem:anchor}, if $K<W_{\mathrm{tot}}$ then the trivial cut $\{r\}$
is infeasible (it has $D(\cut(\{r\}))=W_{\mathrm{tot}}>K$), and every other cut
$X$ with $r\notin X$ corresponds bijectively to a nonempty $S\subsetneq V$ with
\[
  \frac{c_{G'}(\cut(X))}{D(\cut(X))}
  =
  \frac{c_G(\cut(S))}{w_G(S,V)}
  = h(S).
\]
Hence
\[
  \OPT_K(G') = \min_{\substack{S\subsetneq V,\, S\neq\emptyset\\ w_G(S,V)\le K}}
               h(S).
\]
In particular, $\OPT_{K^*}(G')=\psi$ where $K^*=w_G(S^*,V)$ and $S^*$ is any
minimizer of $h$.

\section{Räcke's cut-sparsifier decomposition}

\subsection{Decomposition trees}

A \emph{decomposition tree} for a graph $H=(V',E',c)$ is a rooted tree
$T=(V_T,E_T)$ whose leaves are in bijection with $V'$, together with:
\begin{enumerate}
  \item a node map $m_V:V_T\to V'$ that is bijective on leaves,
  \item an edge map $m_E:E_T\to E'^*$ sending each tree edge to a path in $H$
  between the mapped endpoints,
  \item induced maps $m'_V:V'\to V_T$ (to leaves) and $m'_E:E'\to E_T^*$
  (the unique tree path between mapped endpoints).
\end{enumerate}
For a tree edge $e_t=(u_t,v_t)$, let $V_{u_t},V_{v_t}$ be the leaf partition
induced by deleting $e_t$. Its capacity is
\[
  c_T(e_t):=\sum_{\substack{u\in V_{u_t}\\ v\in V_{v_t}}} c(u,v).
\]

\subsection{Exact Räcke statements used below}

\begin{theorem}[Räcke (STOC 2008, Theorem 2)]
Let $f$ be any multicommodity flow in $H$ with congestion $C_H$.
For any decomposition tree $T$ of $H$, the mapped flow $m'(f)$ in $T$
has congestion $C_T\le C_H$.
\end{theorem}

\begin{theorem}[Räcke (STOC 2008, Theorem 4)]
There exists a polynomial-time computable convex combination of decomposition
trees $(T_i,\lambda_i)$, $\sum_i\lambda_i=1$, and a value
$\beta=O(\log n')$ such that for any family of multicommodity flows
$f_i$ with congestion at most $1$ in $T_i$, the mapped flow
\[
  f:=\sum_i \lambda_i m(f_i)
\]
has congestion at most $\beta$ in $H$.
\end{theorem}

\subsection{Dual cut formulation and proof}

\begin{theorem}[Dual cut bound from Räcke]\label{thm:racke}
Let $H=(V',E',c)$ and let $(T_i,\lambda_i)$ be as in Räcke's Theorem~4.
Define $\mu$ by $\Pr_{T\sim\mu}[T=T_i]=\lambda_i$ and set
$\rho:=\beta=O(\log n')$. Then for every cut $U\subsetneq V'$:
\begin{enumerate}
  \item[(i)] $c_H(\cut(U))\le c_T(\cut(U))$ for every $T$ in the support of $\mu$.
  \item[(ii)] $\Etree[c_T(\cut(U))]\le \rho\,c_H(\cut(U))$.
\end{enumerate}
\end{theorem}

\begin{proof}
Fix a cut $U\subsetneq V'$.

For (i), let $F^*=c_H(\cut(U))$. By max-flow/min-cut in $H$, there is a
single-commodity flow from $U$ to $\bar U$ of value $F^*$ and congestion $1$.
Apply Räcke's Theorem~2 to map this flow to any tree $T$ in support. The mapped
flow in $T$ still has congestion at most $1$, hence can use at most the capacity
of the tree cut $\cut_T(U)$. Therefore
\[
  F^*=c_H(\cut(U))\le c_T(\cut(U)).
\]

For (ii), for each tree $T_i$ choose a maximum $U$--$\bar U$ flow $f_i$ in $T_i$.
Since $T_i$ is a tree and capacities are $c_{T_i}$, max-flow/min-cut gives
\[
  \mathrm{val}(f_i)=c_{T_i}(\cut(U)),
\]
and $f_i$ has congestion at most $1$ in $T_i$.
Apply Räcke's Theorem~4 to the family $(f_i)$: the mapped convex combination
\[
  f:=\sum_i \lambda_i m(f_i)
\]
has congestion at most $\beta$ in $H$. Its value is
\[
  F:=\sum_i \lambda_i\,\mathrm{val}(f_i)
   =\sum_i\lambda_i c_{T_i}(\cut(U))
   =\Etree[c_T(\cut(U))].
\]
Any $U$--$\bar U$ flow of value $F$ in $H$ must have congestion at least
$F/c_H(\cut(U))$ (again by max-flow/min-cut), so
\[
  \frac{F}{c_H(\cut(U))}\le \beta.
\]
Substituting $F=\Etree[c_T(\cut(U))]$ gives
\[
  \Etree[c_T(\cut(U))]\le \beta\,c_H(\cut(U))=\rho\,c_H(\cut(U)).
\]
This proves (ii).
\end{proof}




\section{Approximation for a fixed bound $K$}

\begin{theorem}[Approximation for fixed $K$]\label{thm:fixed-K}
  Apply Theorem~\ref{thm:racke} to $G'$, obtaining distribution $\mu$.
  There exists a tree $T$ in the support of $\mu$ such that:
  \begin{enumerate}
    \item[(a)] Every feasible cut for $\OPT_K(G')$ is also feasible in $T$
      (since feasibility is determined solely by $D$, which is unchanged).
    \item[(b)] $\OPT_K(T)\le \rho\,\OPT_K(G')$.
    \item[(c)] If $X_T$ is an optimal feasible cut of $\OPT_K(T)$, then
      $c_{G'}(\cut(X_T))/D(\cut(X_T))\le \rho\,\OPT_K(G')$.
  \end{enumerate}
\end{theorem}

\begin{proof}
  Let $X^*$ be an optimal cut for $\OPT_K(G')$.
  Since the demand function $D$ is the same in $G'$ and in every tree $T$,
  the cut $X^*$ satisfies $0<D(\cut(X^*))\le K$, hence it is feasible in every tree.
  Therefore:
  \[
    \Etree\!\left[\OPT_K(T)\right]
    \;\le\;
    \Etree\!\left[\frac{c_T(\cut(X^*))}{D(\cut(X^*))}\right]
    \;=\;
    \frac{1}{D(\cut(X^*))}\,\Etree\!\left[c_T(\cut(X^*))\right]
    \;\le\;
    \frac{\rho\,c_{G'}(\cut(X^*))}{D(\cut(X^*))}
    \;=\;
    \rho\,\OPT_K(G').
  \]
  By averaging, some tree $T^{\dagger}$ in the support satisfies
  $\OPT_K(T^{\dagger})\le \rho\,\OPT_K(G')$.  This proves~(b).

  For part~(c), let $X_T$ be optimal for $\OPT_K(T^{\dagger})$.
  By Theorem~\ref{thm:racke}(i), $c_{G'}(\cut(X_T))\le c_{T^{\dagger}}(\cut(X_T))$,
  so:
  \[
    \frac{c_{G'}(\cut(X_T))}{D(\cut(X_T))}
    \;\le\;
    \frac{c_{T^{\dagger}}(\cut(X_T))}{D(\cut(X_T))}
    \;=\;
    \OPT_K(T^{\dagger})
    \;\le\;
    \rho\,\OPT_K(G').
    \qedhere
  \]
\end{proof}

\section{Solving the bounded-demand problem on a tree}

\begin{lemma}[Tree decomposition of cuts]\label{lem:tree-cut}
  Root $T$ at $r$. For each edge $e\in T$, let $A_e\subseteq V'$ be the
  subtree component containing the child endpoint (i.e., not containing $r$).
  Define
  \[
    C_e = c_T(e), \qquad D_e = D(\cut_T(A_e)).
  \]
  For any $S\subseteq V\setminus\{r\}$ with $r\notin S$, the path-cut structure of the
  tree gives
  \[
    c_T(\cut(S)) = \sum_{e\in F(S)} C_e,
    \qquad
    D(\cut(S)) = \sum_{e\in F(S)} D_e,
  \]
  where $F(S)\subseteq E(T)$ is the set of tree edges on the ``frontier'' of $S$:
  those edges $e$ such that $A_e$ contributes to $S$ but $A_{e^+}$ does not
  (formally, $A_e\cap S\neq\emptyset$ and $S\nsupseteq A_e$, or $e$ is on the
  tree-path from some element of $S$ to $r$ and crosses the cut).
\end{lemma}

For a general $S$, the set $F(S)$ can consist of multiple edges.
The key observation is that any such cut is a \emph{weighted average of
per-edge ratios}.

\begin{observation}[Weighted average structure]\label{obs:average}
  For any $S$ with $F(S)\neq\emptyset$ and $D(\cut(S))>0$:
  \[
    \frac{c_T(\cut(S))}{D(\cut(S))}
    \;=\;
    \sum_{e\in F(S)} \lambda_e \cdot \frac{C_e}{D_e},
    \qquad
    \lambda_e = \frac{D_e}{D(\cut(S))}\ge 0,
    \quad
    \sum_{e\in F(S)}\lambda_e = 1.
  \]
\end{observation}

\begin{proof}
  Direct calculation:
  $\dfrac{c_T(\cut(S))}{D(\cut(S))}
  =\dfrac{\sum_e C_e}{\sum_e D_e}
  =\sum_e \dfrac{D_e}{\sum_f D_f}\cdot\dfrac{C_e}{D_e}$.
\end{proof}

\begin{corollary}[Tree optimum is attained by a single edge]\label{cor:single-edge}
  \[
    \OPT_K(T)
    \;=\;
    \min_{\substack{e\in T \\ 0<D_e\le K}} \frac{C_e}{D_e}.
  \]
\end{corollary}

\begin{proof}
  ``$\le$'': For any single edge $e$ with $0<D_e\le K$, the cut $A_e$ is feasible
  for $\OPT_K(T)$ and has ratio $C_e/D_e$.

  ``$\ge$'': Let $S$ be any feasible cut with $D(\cut(S))\le K$.
  Since $D(\cut(S))=\sum_{e\in F(S)}D_e\le K$ and all $D_e\ge 0$, we have
  $D_e\le D(\cut(S))\le K$ for every $e\in F(S)$.
  Also $D_e>0$ whenever $e\in F(S)$ (otherwise the edge contributes zero to both
  numerator and denominator and can be dropped from $F(S)$).
  By Observation~\ref{obs:average}, the ratio of $S$ is a convex combination of
  $\{C_e/D_e: e\in F(S)\}$, so
  \[
    \frac{c_T(\cut(S))}{D(\cut(S))}
    \;\ge\;
    \min_{e\in F(S)} \frac{C_e}{D_e}
    \;\ge\;
    \min_{\substack{e\in T\\0<D_e\le K}} \frac{C_e}{D_e}.
    \qedhere
  \]
\end{proof}

\begin{corollary}
  $\OPT_K(T)$ can be computed in $O(n)$ time by a single scan over the edges of $T$.
\end{corollary}

\section{Handling the unknown $K^*$ by doubling}

The preceding analysis works for a fixed $K$.  In practice $K^*=w_G(S^*,V)$ is
unknown.  We handle this by a standard doubling argument.

\begin{theorem}[Guessing step]\label{thm:guessing}
  Let $\psi=\min_{S}h(S)$ be the true optimum and let $S^*$ be a minimiser with
  $K^*=w_G(S^*,V)$.  Then $K^*\in\{1,2,4,\dots,2^{\lfloor\log_2 W_{\mathrm{tot}}\rfloor}\}$
  if all $\pi(u)$ are integers (or otherwise $K^*\in(W_{\mathrm{tot}}/n,\,W_{\mathrm{tot}}]$
  and a geometric grid suffices).  In either case there is a guess $\hat K$ with
  $K^*\le \hat K\le 2K^*$.
\end{theorem}

\begin{proof}
  Take the smallest power of 2 that is $\ge K^*$.  This gives $\hat K\le 2K^*$.
\end{proof}

\begin{corollary}[Approximation ratio with guessing]\label{cor:approx}
  Using the guess $\hat K$, we have
  $\OPT_{\hat K}(G')\le \psi\le \OPT_{\hat K}(G')$,\footnote{More precisely:
  since $K^*\le\hat K$, the optimum $S^*$ is feasible for $\OPT_{\hat K}(G')$,
  giving $\OPT_{\hat K}(G')\le h(S^*)=\psi$; the other inequality holds because
  every feasible cut for $\hat K$ has $w_G(S,V)\le\hat K$ and has ratio $\ge\psi$.
  Hence $\OPT_{\hat K}(G')=\psi$.} so the algorithm returns a cut $S$ with
  \[
    h(S)\;\le\;\rho\,\psi
    \;=\;
    O(\log n)\cdot\psi.
  \]
\end{corollary}

More precisely, we need not guess: since $\OPT_{\hat K}(G')=\psi$ for the
\emph{correct} scale of $K^*$, we simply run the algorithm for each of the
$O(\log W_{\mathrm{tot}})$ values of $\hat K$ and return the cut with smallest $h(S)$.

\section{Full algorithm}

\begin{algorithm}[H]
\caption{$O(\log n)$-approximation for minimum one-sided quasi-expansion}
\begin{algorithmic}[1]
\Require{Graph $G=(V,E,c,w)$.}
\Ensure{A cut $S^{\dagger}$ with $h(S^{\dagger})\le O(\log n)\cdot\psi$.}
\State Compute $\pi(u)=\sum_v w(u,v)$ for all $u\in V$ and $W_{\mathrm{tot}}=\sum_u\pi(u)$.
\State Build $G'=(V\cup\{r\},E,c,D)$ with $D(u,r)=\pi(u)$.
\State Compute Räcke decomposition $\mu$ over trees of $G'$, with factor $\rho=O(\log n)$.
\State $\text{best}\leftarrow +\infty$; $S^{\dagger}\leftarrow\emptyset$.
\For{each $\hat K\in\{1,2,4,8,\dots\}\cap(0,W_{\mathrm{tot}})$}
  \For{each tree $T$ in the support of $\mu$ (or the single ``winner'' tree in expectation)}
    \State Root $T$ at $r$; compute $\{C_e,D_e\}$ for all edges by a DFS.
    \State Find edge $e^*=\arg\min\{C_e/D_e : 0<D_e\le\hat K\}$.
    \State Let $S \leftarrow A_{e^*}\cap V$ (the subtree at $e^*$, restricted to $V$).
    \If{$h(S)<\text{best}$}
      \State $\text{best}\leftarrow h(S)$; $S^{\dagger}\leftarrow S$.
    \EndIf
  \EndFor
\EndFor
\State \Return $S^{\dagger}$.
\end{algorithmic}
\end{algorithm}

\begin{theorem}[Main result]\label{thm:main}
  The algorithm above runs in polynomial time and returns a cut $S^{\dagger}$ with
  \[
    h(S^{\dagger})\;\le\;O(\log n)\cdot\min_{\emptyset\neq S\subsetneq V}h(S).
  \]
\end{theorem}

\begin{proof}
  Correctness follows from Theorem~\ref{thm:fixed-K}, Corollary~\ref{cor:single-edge},
  and Corollary~\ref{cor:approx}:
  for the iteration with $\hat K\ge K^*$ and $\hat K\le 2K^*$, there exists a tree
  $T^{\dagger}$ in the support of $\mu$ such that $\OPT_{\hat K}(T^{\dagger})\le\rho\,\psi$.
  Corollary~\ref{cor:single-edge} shows this optimum is achieved by a single edge $e^*$,
  yielding $S = A_{e^*}\cap V$ with
  \[
    h(S)
    =\frac{c_{G'}(\cut(X_T))}{D(\cut(X_T))}
    \;\le\;
    \frac{c_{T^{\dagger}}(e^*)}{D_{e^*}}
    =\OPT_{\hat K}(T^{\dagger})
    \;\le\;
    \rho\,\psi.
  \]
  The first inequality uses Theorem~\ref{thm:racke}(i).

  Polynomial runtime: Räcke decomposition runs in polynomial time; iterating over
  $O(\log W_{\mathrm{tot}})=O(\log(n\cdot c_{\max}))=\text{poly}(n)$ values of $\hat K$
  and scanning $O(n)$ tree edges each time is polynomial.
\end{proof}

\end{document}
